
%%%%%%%%%%%%%%%%%%%%%%%%%%%%%%%%%%%%%%%%%%%%%%%%%%%%%%%%%%%%%%%%%%%%%%%%%%%%%%

%\section{Polynomial running time}
%\label{sec: polynomial running time}
%
%%%%%%%%%%%%%%%%%%%%%%%%%%%%%%%%%%%%%%%%%%%%%%%%%%%%%%%%%%%%%%%%%%%%%%%%%%%%%%

%\section{Polynomial running time}
%\label{sec: polynomial running time}
%
%%%%%%%%%%%%%%%%%%%%%%%%%%%%%%%%%%%%%%%%%%%%%%%%%%%%%%%%%%%%%%%%%%%%%%%%%%%%%%

%\section{Polynomial running time}
%\label{sec: polynomial running time}
%
%%%%%%%%%%%%%%%%%%%%%%%%%%%%%%%%%%%%%%%%%%%%%%%%%%%%%%%%%%%%%%%%%%%%%%%%%%%%%%

%\section{Polynomial running time}
%\label{sec: polynomial running time}
%\input{7_polynomial_running_time.tex}

%%%%%%%%
	
	In this section we address Theorem~\ref{thm: miscibility characterization}(c), namely
	the claim that in polynomial time we can test 
	whether a given configuration $C$ is perfectly mixable and, if so, we can
	compute a polynomial-size perfect-mixing graph for $C$ --- also in polynomial time.
	
	Let $C$ be the input configuration. To test whether $C$ is perfectly mixable we just
	need to test whether it satisfies Condition~{\MixCond}.
	That this can be done in polynomial
	time follows directly from Corollary~\ref{cor: mc prime power factors} 
	in Section~\ref{sec: some auxiliary lemmas}.
	To justify this, recall that
	the input size is $s(C) = \sum_{c\in C} \log(|c|+2)\ge n$.
	Thus the factoring of $n$ can be computed in time polynomial in the input size.
	As $n$ has at most $\log n$ distinct odd prime factors and
	each such prime factor has at most $\log{c_{max}}$ powers that are no bigger than $c_{max}$,
	the total number of $b$'s that need to be considered 
	is at most $\log{n}\log{c_{max}}$, which is polynomial in $s(C)$.
	Putting it all together (see the pseudo-code in Algorithm~\ref{alg:perfect mixability}), 
	we obtain that testing perfect-mixability can be indeed accomplished in polynomial time, 
	thus proving the first part of Theorem~\ref{thm: miscibility characterization}(c).
	
	
	\begin{algorithm}[H]
		\caption{PerfectMixabilityTesting($C$)}
		\label{alg:perfect mixability}
		\begin{algorithmic}[1]
			\State{$n\assign\nofdroplets{C}$}
			\State{$\mu\assign \average(C)$}
			\State{$c_{max}\assign$ maximum absolute concentration in $C$}
			\State{$P\assign$ powers of odd prime factors of $n$ that are at most $c_{max}$}
			\ForAll{$p\in P$}
			\If{$C$ is $p$-congruent but $C\cup\braced{\mu}$ is not}
			\State{return false}
			\EndIf
			\EndFor
			\State{return true}
		\end{algorithmic}
	\end{algorithm}


	Now, assume that $C$ is perfectly mixable.
	To prove the second part of Theorem~\ref{thm: miscibility characterization}(c),
	namely that a polynomial-size perfect-mixing graph for $C$ can be constructed in polynomial time,
	the idea is to follow the construction from the proof in Section~\ref{sec: polynomial proof} 
	for Theorem~\ref{thm: miscibility characterization}(b).
	In this proof, at each step we choose a mixing pair in $E_\pi\subseteq E$, for $\pi\in\braced{even,odd}$,
	that is $(\lambda)$-safe (for the appropriate invariant $\lambda\in\braced{I,I'}$), 
	and sufficiently far apart (with respect to either 
	the entire configuration or a subset of the same), until $E$ becomes near-final.
	The former conditions can be checked in polynomial time 
	and there are only quadratically many pairs to try. 
	
	The only remaining obstacle is that, the way the argument is presented above, we would
	have verify at each step whether $E$ is near-final and, if it is, to
	find its near-final partition.  This is not possible in polynomial time (see the remarks below).
	We now explain how to circumvent this problem.
	For $n<22$, then this is simple, since $n$ is a constant, so we focus on the case when $n\ge 22$.
    For our construction in Section~\ref{sec: polynomial - arbitrary n},
	we can just continuously mix $(\lambda)$-safe pairs until $E$ becomes $(\lambda)$-mixed.
	When this happens, as $E$ satisfies Invariant~($\lambda$),
	it holds that there is a near-final pair in $E$.
	Furthermore, as explained in Case~1.1 of the proof of Lemma~\ref{lem: farm mix in E odd overlapping E even},
	this $E$ has the form $E=\braced{f_1:c_1,f_2:c_2,f_3:c_3}$ 
	with $f_1\leq 2$, $f_2=2$ and $\half(c_1+c_2)=c_3$,
	so $(c_1,c_2)$ is a near-final pair.
	In this case, computing the corresponding ``near-final''
	partition of $E$ in polynomial-time is easy:
	if $f_1=1$, then mixing $(c_1,c_2)$ produces $E$ satisfying $\nofconcentrations{E}=2$
	and the near-final partition is given in Lemma~\ref{lem: mixability m = 2}.
	Else $f_1=2$ and the near-final partition involves sets of the form 
	$\braced{c_1,c_2}$ and $\braced{c_3}$, where $c_3=\average(E)$.

%%%%%%%%%

\emph{Remarks.}
The last two paragraphs bring up a question of whether checking the near-final property can be
done in polynomial time. While this is not essential to our algorithm, we address it here
for the sake of interested readers.

Let $\nofdroplets{E} = n$. Observe first that if $E$ is near-final then it has a
near-final partition where all sets in the partition have different cardinalities.
The reason is simple: if two sets have the same cardinality, we can merge them into
one set, retaining the average value and the property that all cardinalities are powers of $2$.
As a consequence of this, the cardinalities of the sets in $E$'s partition are
uniquely determined by $n$, namely these are the powers of $2$ that appear in the binary
representation of $n$. Let $\omega = \mysum(E) = n\cdot\hatmu$
be the sum of all concentrations in $E$, and let
$k$ be some cardinality (a power of $2$) in this partition of $E$. 
A subset $F\subseteq E$ of cardinality $k$ has $\average(F) = \hatmu$ if and only if
$\mysum(F) = \omega k/n$. 

The above properties lead to an $\NP$-completeness proof.
The idea is this. Consider a restricted variant of the problem where
$n = 3\cdot 2^l$, for some $l$. Using the notation from the paragraph above, we need to
determine whether $E$ has a subset $F$ of cardinality $k = n/3$ and
total sum $\mysum(F) = \omega/3$. This is a variant of the $\Partition$
problem that can easily be shown to be $\NP$-complete. (Details are left to the reader.)











%%%%%%%%
	
	In this section we address Theorem~\ref{thm: miscibility characterization}(c), namely
	the claim that in polynomial time we can test 
	whether a given configuration $C$ is perfectly mixable and, if so, we can
	compute a polynomial-size perfect-mixing graph for $C$ --- also in polynomial time.
	
	Let $C$ be the input configuration. To test whether $C$ is perfectly mixable we just
	need to test whether it satisfies Condition~{\MixCond}.
	That this can be done in polynomial
	time follows directly from Corollary~\ref{cor: mc prime power factors} 
	in Section~\ref{sec: some auxiliary lemmas}.
	To justify this, recall that
	the input size is $s(C) = \sum_{c\in C} \log(|c|+2)\ge n$.
	Thus the factoring of $n$ can be computed in time polynomial in the input size.
	As $n$ has at most $\log n$ distinct odd prime factors and
	each such prime factor has at most $\log{c_{max}}$ powers that are no bigger than $c_{max}$,
	the total number of $b$'s that need to be considered 
	is at most $\log{n}\log{c_{max}}$, which is polynomial in $s(C)$.
	Putting it all together (see the pseudo-code in Algorithm~\ref{alg:perfect mixability}), 
	we obtain that testing perfect-mixability can be indeed accomplished in polynomial time, 
	thus proving the first part of Theorem~\ref{thm: miscibility characterization}(c).
	
	
	\begin{algorithm}[H]
		\caption{PerfectMixabilityTesting($C$)}
		\label{alg:perfect mixability}
		\begin{algorithmic}[1]
			\State{$n\assign\nofdroplets{C}$}
			\State{$\mu\assign \average(C)$}
			\State{$c_{max}\assign$ maximum absolute concentration in $C$}
			\State{$P\assign$ powers of odd prime factors of $n$ that are at most $c_{max}$}
			\ForAll{$p\in P$}
			\If{$C$ is $p$-congruent but $C\cup\braced{\mu}$ is not}
			\State{return false}
			\EndIf
			\EndFor
			\State{return true}
		\end{algorithmic}
	\end{algorithm}


	Now, assume that $C$ is perfectly mixable.
	To prove the second part of Theorem~\ref{thm: miscibility characterization}(c),
	namely that a polynomial-size perfect-mixing graph for $C$ can be constructed in polynomial time,
	the idea is to follow the construction from the proof in Section~\ref{sec: polynomial proof} 
	for Theorem~\ref{thm: miscibility characterization}(b).
	In this proof, at each step we choose a mixing pair in $E_\pi\subseteq E$, for $\pi\in\braced{even,odd}$,
	that is $(\lambda)$-safe (for the appropriate invariant $\lambda\in\braced{I,I'}$), 
	and sufficiently far apart (with respect to either 
	the entire configuration or a subset of the same), until $E$ becomes near-final.
	The former conditions can be checked in polynomial time 
	and there are only quadratically many pairs to try. 
	
	The only remaining obstacle is that, the way the argument is presented above, we would
	have verify at each step whether $E$ is near-final and, if it is, to
	find its near-final partition.  This is not possible in polynomial time (see the remarks below).
	We now explain how to circumvent this problem.
	For $n<22$, then this is simple, since $n$ is a constant, so we focus on the case when $n\ge 22$.
    For our construction in Section~\ref{sec: polynomial - arbitrary n},
	we can just continuously mix $(\lambda)$-safe pairs until $E$ becomes $(\lambda)$-mixed.
	When this happens, as $E$ satisfies Invariant~($\lambda$),
	it holds that there is a near-final pair in $E$.
	Furthermore, as explained in Case~1.1 of the proof of Lemma~\ref{lem: farm mix in E odd overlapping E even},
	this $E$ has the form $E=\braced{f_1:c_1,f_2:c_2,f_3:c_3}$ 
	with $f_1\leq 2$, $f_2=2$ and $\half(c_1+c_2)=c_3$,
	so $(c_1,c_2)$ is a near-final pair.
	In this case, computing the corresponding ``near-final''
	partition of $E$ in polynomial-time is easy:
	if $f_1=1$, then mixing $(c_1,c_2)$ produces $E$ satisfying $\nofconcentrations{E}=2$
	and the near-final partition is given in Lemma~\ref{lem: mixability m = 2}.
	Else $f_1=2$ and the near-final partition involves sets of the form 
	$\braced{c_1,c_2}$ and $\braced{c_3}$, where $c_3=\average(E)$.

%%%%%%%%%

\emph{Remarks.}
The last two paragraphs bring up a question of whether checking the near-final property can be
done in polynomial time. While this is not essential to our algorithm, we address it here
for the sake of interested readers.

Let $\nofdroplets{E} = n$. Observe first that if $E$ is near-final then it has a
near-final partition where all sets in the partition have different cardinalities.
The reason is simple: if two sets have the same cardinality, we can merge them into
one set, retaining the average value and the property that all cardinalities are powers of $2$.
As a consequence of this, the cardinalities of the sets in $E$'s partition are
uniquely determined by $n$, namely these are the powers of $2$ that appear in the binary
representation of $n$. Let $\omega = \mysum(E) = n\cdot\hatmu$
be the sum of all concentrations in $E$, and let
$k$ be some cardinality (a power of $2$) in this partition of $E$. 
A subset $F\subseteq E$ of cardinality $k$ has $\average(F) = \hatmu$ if and only if
$\mysum(F) = \omega k/n$. 

The above properties lead to an $\NP$-completeness proof.
The idea is this. Consider a restricted variant of the problem where
$n = 3\cdot 2^l$, for some $l$. Using the notation from the paragraph above, we need to
determine whether $E$ has a subset $F$ of cardinality $k = n/3$ and
total sum $\mysum(F) = \omega/3$. This is a variant of the $\Partition$
problem that can easily be shown to be $\NP$-complete. (Details are left to the reader.)











%%%%%%%%
	
	In this section we address Theorem~\ref{thm: miscibility characterization}(c), namely
	the claim that in polynomial time we can test 
	whether a given configuration $C$ is perfectly mixable and, if so, we can
	compute a polynomial-size perfect-mixing graph for $C$ --- also in polynomial time.
	
	Let $C$ be the input configuration. To test whether $C$ is perfectly mixable we just
	need to test whether it satisfies Condition~{\MixCond}.
	That this can be done in polynomial
	time follows directly from Corollary~\ref{cor: mc prime power factors} 
	in Section~\ref{sec: some auxiliary lemmas}.
	To justify this, recall that
	the input size is $s(C) = \sum_{c\in C} \log(|c|+2)\ge n$.
	Thus the factoring of $n$ can be computed in time polynomial in the input size.
	As $n$ has at most $\log n$ distinct odd prime factors and
	each such prime factor has at most $\log{c_{max}}$ powers that are no bigger than $c_{max}$,
	the total number of $b$'s that need to be considered 
	is at most $\log{n}\log{c_{max}}$, which is polynomial in $s(C)$.
	Putting it all together (see the pseudo-code in Algorithm~\ref{alg:perfect mixability}), 
	we obtain that testing perfect-mixability can be indeed accomplished in polynomial time, 
	thus proving the first part of Theorem~\ref{thm: miscibility characterization}(c).
	
	
	\begin{algorithm}[H]
		\caption{PerfectMixabilityTesting($C$)}
		\label{alg:perfect mixability}
		\begin{algorithmic}[1]
			\State{$n\assign\nofdroplets{C}$}
			\State{$\mu\assign \average(C)$}
			\State{$c_{max}\assign$ maximum absolute concentration in $C$}
			\State{$P\assign$ powers of odd prime factors of $n$ that are at most $c_{max}$}
			\ForAll{$p\in P$}
			\If{$C$ is $p$-congruent but $C\cup\braced{\mu}$ is not}
			\State{return false}
			\EndIf
			\EndFor
			\State{return true}
		\end{algorithmic}
	\end{algorithm}


	Now, assume that $C$ is perfectly mixable.
	To prove the second part of Theorem~\ref{thm: miscibility characterization}(c),
	namely that a polynomial-size perfect-mixing graph for $C$ can be constructed in polynomial time,
	the idea is to follow the construction from the proof in Section~\ref{sec: polynomial proof} 
	for Theorem~\ref{thm: miscibility characterization}(b).
	In this proof, at each step we choose a mixing pair in $E_\pi\subseteq E$, for $\pi\in\braced{even,odd}$,
	that is $(\lambda)$-safe (for the appropriate invariant $\lambda\in\braced{I,I'}$), 
	and sufficiently far apart (with respect to either 
	the entire configuration or a subset of the same), until $E$ becomes near-final.
	The former conditions can be checked in polynomial time 
	and there are only quadratically many pairs to try. 
	
	The only remaining obstacle is that, the way the argument is presented above, we would
	have verify at each step whether $E$ is near-final and, if it is, to
	find its near-final partition.  This is not possible in polynomial time (see the remarks below).
	We now explain how to circumvent this problem.
	For $n<22$, then this is simple, since $n$ is a constant, so we focus on the case when $n\ge 22$.
    For our construction in Section~\ref{sec: polynomial - arbitrary n},
	we can just continuously mix $(\lambda)$-safe pairs until $E$ becomes $(\lambda)$-mixed.
	When this happens, as $E$ satisfies Invariant~($\lambda$),
	it holds that there is a near-final pair in $E$.
	Furthermore, as explained in Case~1.1 of the proof of Lemma~\ref{lem: farm mix in E odd overlapping E even},
	this $E$ has the form $E=\braced{f_1:c_1,f_2:c_2,f_3:c_3}$ 
	with $f_1\leq 2$, $f_2=2$ and $\half(c_1+c_2)=c_3$,
	so $(c_1,c_2)$ is a near-final pair.
	In this case, computing the corresponding ``near-final''
	partition of $E$ in polynomial-time is easy:
	if $f_1=1$, then mixing $(c_1,c_2)$ produces $E$ satisfying $\nofconcentrations{E}=2$
	and the near-final partition is given in Lemma~\ref{lem: mixability m = 2}.
	Else $f_1=2$ and the near-final partition involves sets of the form 
	$\braced{c_1,c_2}$ and $\braced{c_3}$, where $c_3=\average(E)$.

%%%%%%%%%

\emph{Remarks.}
The last two paragraphs bring up a question of whether checking the near-final property can be
done in polynomial time. While this is not essential to our algorithm, we address it here
for the sake of interested readers.

Let $\nofdroplets{E} = n$. Observe first that if $E$ is near-final then it has a
near-final partition where all sets in the partition have different cardinalities.
The reason is simple: if two sets have the same cardinality, we can merge them into
one set, retaining the average value and the property that all cardinalities are powers of $2$.
As a consequence of this, the cardinalities of the sets in $E$'s partition are
uniquely determined by $n$, namely these are the powers of $2$ that appear in the binary
representation of $n$. Let $\omega = \mysum(E) = n\cdot\hatmu$
be the sum of all concentrations in $E$, and let
$k$ be some cardinality (a power of $2$) in this partition of $E$. 
A subset $F\subseteq E$ of cardinality $k$ has $\average(F) = \hatmu$ if and only if
$\mysum(F) = \omega k/n$. 

The above properties lead to an $\NP$-completeness proof.
The idea is this. Consider a restricted variant of the problem where
$n = 3\cdot 2^l$, for some $l$. Using the notation from the paragraph above, we need to
determine whether $E$ has a subset $F$ of cardinality $k = n/3$ and
total sum $\mysum(F) = \omega/3$. This is a variant of the $\Partition$
problem that can easily be shown to be $\NP$-complete. (Details are left to the reader.)











%%%%%%%%
	
	In this section we address Theorem~\ref{thm: miscibility characterization}(c), namely
	the claim that in polynomial time we can test 
	whether a given configuration $C$ is perfectly mixable and, if so, we can
	compute a polynomial-size perfect-mixing graph for $C$ --- also in polynomial time.
	
	Let $C$ be the input configuration. To test whether $C$ is perfectly mixable we just
	need to test whether it satisfies Condition~{\MixCond}.
	That this can be done in polynomial
	time follows directly from Corollary~\ref{cor: mc prime power factors} 
	in Section~\ref{sec: some auxiliary lemmas}.
	To justify this, recall that
	the input size is $s(C) = \sum_{c\in C} \log(|c|+2)\ge n$.
	Thus the factoring of $n$ can be computed in time polynomial in the input size.
	As $n$ has at most $\log n$ distinct odd prime factors and
	each such prime factor has at most $\log{c_{max}}$ powers that are no bigger than $c_{max}$,
	the total number of $b$'s that need to be considered 
	is at most $\log{n}\log{c_{max}}$, which is polynomial in $s(C)$.
	Putting it all together (see the pseudo-code in Algorithm~\ref{alg:perfect mixability}), 
	we obtain that testing perfect-mixability can be indeed accomplished in polynomial time, 
	thus proving the first part of Theorem~\ref{thm: miscibility characterization}(c).
	
	
	\begin{algorithm}[H]
		\caption{PerfectMixabilityTesting($C$)}
		\label{alg:perfect mixability}
		\begin{algorithmic}[1]
			\State{$n\assign\nofdroplets{C}$}
			\State{$\mu\assign \average(C)$}
			\State{$c_{max}\assign$ maximum absolute concentration in $C$}
			\State{$P\assign$ powers of odd prime factors of $n$ that are at most $c_{max}$}
			\ForAll{$p\in P$}
			\If{$C$ is $p$-congruent but $C\cup\braced{\mu}$ is not}
			\State{return false}
			\EndIf
			\EndFor
			\State{return true}
		\end{algorithmic}
	\end{algorithm}


	Now, assume that $C$ is perfectly mixable.
	To prove the second part of Theorem~\ref{thm: miscibility characterization}(c),
	namely that a polynomial-size perfect-mixing graph for $C$ can be constructed in polynomial time,
	the idea is to follow the construction from the proof in Section~\ref{sec: polynomial proof} 
	for Theorem~\ref{thm: miscibility characterization}(b).
	In this proof, at each step we choose a mixing pair in $E_\pi\subseteq E$, for $\pi\in\braced{even,odd}$,
	that is $(\lambda)$-safe (for the appropriate invariant $\lambda\in\braced{I,I'}$), 
	and sufficiently far apart (with respect to either 
	the entire configuration or a subset of the same), until $E$ becomes near-final.
	The former conditions can be checked in polynomial time 
	and there are only quadratically many pairs to try. 
	
	The only remaining obstacle is that, the way the argument is presented above, we would
	have verify at each step whether $E$ is near-final and, if it is, to
	find its near-final partition.  This is not possible in polynomial time (see the remarks below).
	We now explain how to circumvent this problem.
	For $n<22$, then this is simple, since $n$ is a constant, so we focus on the case when $n\ge 22$.
    For our construction in Section~\ref{sec: polynomial - arbitrary n},
	we can just continuously mix $(\lambda)$-safe pairs until $E$ becomes $(\lambda)$-mixed.
	When this happens, as $E$ satisfies Invariant~($\lambda$),
	it holds that there is a near-final pair in $E$.
	Furthermore, as explained in Case~1.1 of the proof of Lemma~\ref{lem: farm mix in E odd overlapping E even},
	this $E$ has the form $E=\braced{f_1:c_1,f_2:c_2,f_3:c_3}$ 
	with $f_1\leq 2$, $f_2=2$ and $\half(c_1+c_2)=c_3$,
	so $(c_1,c_2)$ is a near-final pair.
	In this case, computing the corresponding ``near-final''
	partition of $E$ in polynomial-time is easy:
	if $f_1=1$, then mixing $(c_1,c_2)$ produces $E$ satisfying $\nofconcentrations{E}=2$
	and the near-final partition is given in Lemma~\ref{lem: mixability m = 2}.
	Else $f_1=2$ and the near-final partition involves sets of the form 
	$\braced{c_1,c_2}$ and $\braced{c_3}$, where $c_3=\average(E)$.

%%%%%%%%%

\emph{Remarks.}
The last two paragraphs bring up a question of whether checking the near-final property can be
done in polynomial time. While this is not essential to our algorithm, we address it here
for the sake of interested readers.

Let $\nofdroplets{E} = n$. Observe first that if $E$ is near-final then it has a
near-final partition where all sets in the partition have different cardinalities.
The reason is simple: if two sets have the same cardinality, we can merge them into
one set, retaining the average value and the property that all cardinalities are powers of $2$.
As a consequence of this, the cardinalities of the sets in $E$'s partition are
uniquely determined by $n$, namely these are the powers of $2$ that appear in the binary
representation of $n$. Let $\omega = \mysum(E) = n\cdot\hatmu$
be the sum of all concentrations in $E$, and let
$k$ be some cardinality (a power of $2$) in this partition of $E$. 
A subset $F\subseteq E$ of cardinality $k$ has $\average(F) = \hatmu$ if and only if
$\mysum(F) = \omega k/n$. 

The above properties lead to an $\NP$-completeness proof.
The idea is this. Consider a restricted variant of the problem where
$n = 3\cdot 2^l$, for some $l$. Using the notation from the paragraph above, we need to
determine whether $E$ has a subset $F$ of cardinality $k = n/3$ and
total sum $\mysum(F) = \omega/3$. This is a variant of the $\Partition$
problem that can easily be shown to be $\NP$-complete. (Details are left to the reader.)









